\begin{solution}{15}
\begin{subsolution}
考虑上侧以 $45^{\circ}$ 角入射到细丝外表面的细光束 $AP$ ，它被经细丝外表面反射后成方向平行于 $z$ 轴正向的反射光 $PB$ ，其反向延长线与 $y$ 轴交于 $Q$ 点，如图所示。由几何关系可知，光线 $PB$ 离 $xz$ 平面的距离为
\begin{equation}
O Q = + \frac {\sqrt {2}}{4} d = y _ {+}
\label{eq:1.s15}
\end{equation}
考虑上侧以 $45^{\circ} + \Delta\theta$ （ $\Delta\theta$ 很小）角入射到细丝外表面的细光束 $A_{+}P_{+}$ ，它被经细丝外表面反射后的反射光 $P_{+}B_{+}$ ，其反向延长线与 $y$ 轴交于 $E$ 点，如图所示。由反射定律可知
\begin{equation}
\angle A _ {+} P _ {+} B _ {+} = 2 \times (45 ^ {\circ} + \Delta \theta)
\label{eq:2.s15}
\end{equation}
因而，反射光 $P_{+}B_{+}$ 与 $\mathbf{z}$ 轴的夹角为
\begin{equation}
2 \Delta \theta
\label{eq:3.s15}
\end{equation}
考虑 $\Delta EOP_{+}$ ，由几何关系可知
\begin{equation*}
\angle E P _ {+} O = 45 ^ {\circ} + \Delta \theta = \angle E O P _ {+}
\end{equation*}
于是
\begin{equation}
O E = P _ {+} E
\label{eq:4.s15}
\end{equation}
由于 $OP_{+}=d/2$ ，所以
\begin{equation*}
\frac {d}{4} = O E \cos (45 ^ {\circ} + \Delta \theta)
\end{equation*}
利用
\begin{equation*}
\cos (45 ^ {\circ} + \Delta \theta) = \cos 45 ^ {\circ} \cos \Delta \theta - \sin 45 ^ {\circ} \sin \Delta \theta
\end{equation*}
和
\begin{equation*}
\cos \Delta \theta \sim 1, \sin \Delta \theta \sim \Delta \theta , \text {当} \Delta \theta \sim 0
\end{equation*}
可得
\begin{equation*}
\cos (45 ^ {\circ} + \Delta \theta) = \frac {\sqrt {2}}{2} (1 - \Delta \theta)
\end{equation*}
所以
\begin{equation}
O E = \frac {\sqrt {2} d}{4 (1 - \Delta \theta)} \approx \frac {\sqrt {2} d}{4} (1 + \Delta \theta)
\label{eq:5.s15}
\end{equation}
因而由\eqref{eq:1.s15}、\eqref{eq:5.s15}式得
\begin{equation}
Q E = O E - O Q = \frac {\sqrt {2} d}{4} \Delta \theta
\label{eq:6.s15}
\end{equation}
进而，设 $EB_{+}$ 交 $QB$ 于 $F$ （设其 $\mathbf{z}$ 坐标为 $z_{+}$ ），则在直角三角形 $EQF$ 中有
\begin{equation}
z _ {+} = Q F = \frac {Q E}{\tan (2 \Delta \theta)} = \frac {\sqrt {2}}{4} d \frac {\Delta \theta}{2 \Delta \theta} = \frac {\sqrt {2}}{8} d
\label{eq:7.s15}
\end{equation}
因此，上侧入射角在 $\pi /4$ 附近的入射光线虚像的位置为
\begin{equation}
y _ {+} = \frac {\sqrt {2}}{4} d, z _ {+} = \frac {\sqrt {2}}{8} d
\label{eq:8.s15}
\end{equation}
利用对称性可见，下侧入射角在 $\pi /4$ 附近的入射光线虚像的位置为
\begin{equation}
y _ {-} = - \frac {\sqrt {2}}{4} d, z _ {-} = \frac {\sqrt {2}}{8} d
\label{eq:9.s15}
\end{equation}
\end{subsolution}

\begin{subsolution}
本问题可视为由上、下两侧的反射光束分别形成的虚光源所发出的光线之间的干涉，相当于双缝位置分别在 $y_{+} = \frac{\sqrt{2}}{4} d$ ， $z_{+} = \frac{\sqrt{2}}{8} d$ 和 $y_{-} = -\frac{\sqrt{2}}{4} d$ ， $z_{-} = \frac{\sqrt{2}}{8} d$ 的双缝干涉。双缝屏与观察屏之间的距离为 $D - \frac{\sqrt{2}d}{8} \approx D$ ，双缝间距为 $\frac{\sqrt{2}}{2} d$ 。故条纹间距
\begin{equation}
a = \lambda \frac {D}{(\sqrt {2} / 2) d} = \frac {\sqrt {2} \lambda D}{d}
\label{eq:10.s15}
\end{equation}
由此得，细丝的直径 $d$ 为
\begin{equation}
d = \frac {\sqrt {2} \lambda D}{a}
\label{eq:11.s15}
\end{equation}
\end{subsolution}
\end{solution}